\documentclass[a4paper,12pt]{article}
\usepackage[T2A]{fontenc}
\usepackage[utf8]{inputenc}
\usepackage[english,russian]{babel}
\usepackage{natbib}
\usepackage{hyperref}
\hypersetup{
  colorlinks=true,
  linkcolor=blue,
  citecolor=red,
  urlcolor=cyan
}

\title{Лабораторная работа №6 \\ Управление библиографией в LaTeX}
\author{Сунь Маосин}
\date{2026}

\begin{document}

\maketitle

\section{Цель работы}

Изучение возможностей LaTeX по управлению библиографическими ссылками.

\section{Цитирование с помощью natbib}

\subsection{Текстовое цитирование}

\citet{smith2023} показали, что LaTeX улучшает качество научных работ.

\citet{latex2024} в своей книге подробно описывают все возможности LaTeX.

\subsection{Цитирование в скобках}

Как показано в исследованиях \citep{konference2023}, автоматизация публикаций становится всё более важной.

Подробное описание LaTeX можно найти в \citep{latex2024}.

\subsection{Цитирование с страницами}

\citet[p.~45]{smith2023} утверждают, что...

\citep[см.][глава 2]{latex2024} содержит основы LaTeX.

\subsection{Множественные цитирования}

\citep{smith2023,latex2024,konference2023}

\section{Список литературы}

\bibliographystyle{plainnat}
\bibliography{references}

\end{document}